\begin{maintheorem}[Component count and the smooth general fibre]
\hypertarget{tag-03K3}{}\label{series:II:B}

Let $f:Y\to C$ be a proper holomorphic surjection with connected
fibres from a connected smooth compact complex threefold to a
smooth compact curve. Delete a nonempty finite set $\Delta$
containing all critical values, put $s=|\Delta|$, $g=g(C)$,
let $r$ count the irreducible components over $\Delta$, and set
\[
 t'=\dim_\Q H_1(F,\Q)_{\pi_1(C\setminus\Delta)}
\]
for a smooth fibre $F$. With $b_2(Y)=0$,
\[
 r=2g+s-1+t'-b_1(Y).
\]
Under $H^1(Y,\Z)=H^2(Y,\Z)=0$, the base is $\PP^1$,
$r=s-1+t'$, and the integral coinvariant quotient
$H_1(F,\Z)_{\pi_1(C\setminus\Delta)}$ is free abelian.
For a holomorphic algebraic reduction with $a(Y)=1$ and
$b_1(Y)=b_2(Y)=0$, every smooth fibre satisfies
\[
 c_1(K_F)=0\text{ in }H^2(F,\Q),\qquad
 c_2(F)=0,\qquad \chi(F,\OO_F)=0.
\]
Its very general fibre is a complex two-torus or a primary or
secondary Hopf surface of algebraic dimension zero. Each of these
surviving types occurs in the constructions stated in the body.

\end{maintheorem}
